\section{Policy Enforcement 与 Contextual Security 案例}

讲者用 Progent 类方案展示了策略执行的核心思想：用 DSL 显式描述安全策略，在 Agent 执行过程中做实时 policy checking，并允许基于上下文动态更新策略。相比静态策略，这种方法更适合任务驱动且状态不断变化的 Agent。

\begin{knowledgebox}{Contextual Security 的直觉}
同一个 Agent 在不同阶段应有不同权限。例如任务是 ``总结最近一天未读邮件''，此时允许读邮件元数据，但不应默认允许外发邮件。只有在后续状态显式出现 ``用户确认发送'' 才能临时提升发送权限。
\end{knowledgebox}

\begin{lstlisting}[caption={策略示意：基于上下文限制高风险动作}]
policy "email_summary" {
  allow tool.read_email_metadata
  deny  tool.send_email
}

on state(user_confirmed_send == true) {
  allow tool.send_email(to=resolved_recipient)
}
\end{lstlisting}

\begin{importantbox}{执行时策略比离线对齐更直接}
离线对齐提高了模型整体倾向，执行时策略直接约束当前动作。两者叠加，才能在 ``模型被诱导'' 的情况下仍阻断高危行为。
\end{importantbox}

\subsection{本章小结}
Contextual Security 把权限控制从 ``谁'' 扩展到 ``在什么状态下能做什么''。这是 Agent 安全从静态访问控制走向动态行为控制的关键一步。

